\section{Terza Missione e Trasferimento Tecnologico (es. brevetti)}

\outerlist{

\entrybig
	{\textbf{Responsabile scientifico di progetto}}{Dipartimento di Ingegneria “Enzo Ferrari”}
	{}{Collaborazione con Data Reply S.r.l. (2025, 6 mesi)}
\innerlist{
    \entry{Responsabile scientifico del progetto \textit{“Apprendimento online e Lifelong Learning in sistemi di AI”}.}
    \entry{\textbf{Attività:} sviluppo di una libreria \textit{open-source} in linguaggio Python per la progettazione e il benchmarking di algoritmi di \textit{Continual Learning} in ambito \textit{Artificial Intelligence}.}
}

\entrybig
	{\textbf{Responsabile scientifico di progetto}}{Centro di Ricerca AIRI}
	{}{Collaborazione con SEW-EURODRIVE S.a.s. (2025, 12 mesi)}
\innerlist{
    \entry{Coordinatore scientifico del progetto \textit{“SSL neural network definition for Condition Based Monitoring to be applied on different machines type within a classic packaging line”}.}
    \entry{\textbf{Attività:} progettazione e sviluppo di un sistema di \textit{anomaly detection} non supervisionato per applicazioni di manutenzione predittiva su motoriduttori industriali.}
}
\entrybig
	{\textbf{Responsabile tecnico di progetto}}{Dipartimento di Ingegneria “Enzo Ferrari”}
	{}{Collaborazione con ENI S.p.A. (2025, durata 18 mesi)}
\innerlist{
    \entry{\textbf{Attività:} attività di ricerca, sviluppo e supporto tecnico nei settori della \textit{Intelligent Vision} e della \textit{Multimodal Artificial Intelligence}. Progettazione e sviluppo di un sistema di monitoraggio basato su \textit{Computer Vision} per l’analisi automatica di azioni salienti, includendo tecniche di \textit{multi-person detection and tracking} e \textit{open-vocabulary action recognition}, con integrazione di informazioni documentali multimodali.}
}

\entrybig
	{\textbf{Brevetto per invenzione industriale} [\href{https://patents.google.com/patent/US11672255B2/en}{link}]}{Inventore}
	{}{Ministero dello Sviluppo Economico, Italia (2019)}
\innerlist{
    \entry{Inventore di un brevetto per l’analisi di immagini di carcasse animali mediante tecniche di \textit{Artificial Intelligence} basate su \textit{Convolutional Neural Networks} (CNN), sviluppato in collaborazione con Farm4Trade.}
    \entry{\textit{“Metodo di valutazione di uno stato di salute di un elemento anatomico, relativo dispositivo di valutazione e relativo sistema di valutazione”}.}
}

\entrybig
	{\textbf{Organizzazione e public engagement}}{AImageLab}
	{Notte della Ricerca [\href{https://www.nottedellaricerca.unimore.it/2025/evento.html?e=45}{link}]}{UNIMORE (2024--2025)}
\innerlist{
    \entry{Organizzatore dello stand del laboratorio \textit{AImageLab} in occasione della \textit{Notte della Ricerca} (edizioni 2024 e 2025), evento di divulgazione scientifica rivolto al pubblico generale con demo interattive, poster e presentazioni sulle principali linee di ricerca del laboratorio.}
}

\entrybig
{\textbf{Attività di formazione e didattica rivolta ad aziende, centri di competenza e istituzioni}}{}
{\textcolor{darkgray}{Dal 2018, \textbf{oltre 200 ore} complessive di attività di formazione erogata a personale di aziende.}}{}
\innerlist{
\entry{\textit{Master in Machine Learning and Deep Learning} (Fondazione Democenter): 6 edizioni (2018--2024), 16 ore per edizione.}
\entry{\textit{Python, Deep Learning and Computer Vision} (UNIMORE): 24 ore complessive (edizione 2025).}
\entry{\textit{Machine Learning e Deep Learning – Formazione e trasferimento di know-how} (Academy UNIPOL): 1 edizione (2025), 12 ore.}
\entry{\textit{Corso di Machine Learning e Deep Learning} (BI-REX per SACMI): 1 edizione (2025), 12 ore.}
\entry{\textit{Corso teorico e pratico in Machine Learning e Deep Learning -- Master Teknè} (BI-REX): 3 edizioni (2022, 2024, 2025), 8 ore per edizione.}
\entry{\textit{Intensive Mater AI and ML for Smart Factory} (Experis S.r.l.): 1 edizione (2023--2024), 12 ore.}
\entry{\textit{Corso di formazione teorico e pratico su Computer Vision e Machine Learning} (CNH Industrial Italia S.p.A.): 1 edizione (2023), 8 ore.}
\entry{\textit{School in AI: Deep Learning, Vision and Language for Industry} (AI Academy UNIMORE): 2 edizioni (2022), 4 ore per edizione.}
\entry{\textit{Predictive Analytics and Machine Learning} (IMA S.p.A. Unipersonale): 1 edizione (2022), 6 ore.}
\entry{\textit{Training course on Artificial Intelligence} (Istituto Zooprofilattico Sperimentale dell’Abruzzo e del Molise ``G.~Caporale''): 1 edizione (2018), 8 ore.}
}
}
